\clearpage
\section{Distribution weights and technology levels: M11, M13--M15 (Group B)}

The authoritative V6 definitions are
\[
Y_i=\left[\alpha_i(A_{X,i}X_i)^{1-\rho_i}
 +(1-\alpha_i)(A_{L,i}L_i)^{1-\rho_i}\right]^{1/(1-\rho_i)},
\qquad X_i=\varphi_iH_i .
\]
M11 comprises $\alpha_{US},\alpha_W$. M13 comprises
$A^u_{X,US}=\bar A_{X,US,u}N_{US}^{\xi_u}$ and
$A^u_{L,US}=\bar A_{L,US,u}N_{US}^{\nu_u}$.
M14 comprises $A^b_{X,US}=\bar A_{X,US,b}$ and
$A^b_{L,US}=\bar A_{L,US,b}$ because $\nu_b=0$.
M15 comprises $A_{X,W}=\bar A_{X,W}$ and $A_{L,W}=\bar A_{L,W}$ because $\xi_W=0$.
These six barred coefficients are augmentation scales, not observed aggregate TFP.

\begin{evidence}
Caselli and Coleman (2006) \src{CC06}.
Skill-specific augmentations $A_s,A_u$ in a labor CES inside a capital--labor technology;
equations (2)--(3), pp.~499--502.
&
Baseline skill elasticity $\sigma=1.4$ and capital share $1/3$ are imposed.
Augmentations are jointly recovered from output and the skill premium.
There is no unit-free common level or range for V6's barred coefficients.
&
Cross-country development accounting: 52 countries; output/capital 1988 and
schooling 1985. Baseline skill split is completed primary schooling.
Static country comparison; PPP quantities and wage-premium inputs.
&
Transfer the inversion method, not the source capital share to V6 $\alpha$.
Adjust for V6's intermediate markup, production-only skilled labor and omitted physical
capital. Cross-country technology differences are not time-series elasticities in $N$.
\\[5pt]
Rossi (2022) \src{Rossi22}.
Relative skill efficiency $(A_HQ_H)/(A_LQ_L)$; equations (4)--(8), Table~I
in the July~2021 accepted manuscript, pp.~7--8,31.
&
At imposed $\sigma=1.5$, US relative efficiency is normalized to 1;
Canada is $0.711$, India $0.041$. These are country inversions,
not a confidence interval or V6 RoW scale range.
&
Comparable wage and hours microdata for 12 countries around 2000.
Some tertiary education defines high skill.
Separate migrant evidence distinguishes embodied human capital $Q$ from technology.
Static CES, competitive labor pricing.
&
Useful for relative-productivity measurement. The reported object mixes
technology and worker quality, and lacks V6's research-labor and markup adjustments.
Re-invert the chosen aggregate; averaging these country numbers does not produce
$(\bar A_{X,W},\bar A_{L,W})$.
\\[5pt]
Cantore and Levine (2012) \src{CL12};
Klump, McAdam and Willman (2012) \src{KMW12}.
CES distribution and scale normalization.
&
No independent empirical point estimate for V6 $\alpha$ or augmentation levels.
Distribution coefficients depend on units and the normalization point.
Normalized baseline income shares are different objects from an unrestricted raw CES weight.
&
Methodological macroeconomic analyses; RBC examples and a theory/empirics survey.
Detailed checks use the 2011 Dynare and ECB working-paper versions,
respectively. No common population or sample-based confidence interval.
&
For $\rho\ne1$, fix one convention and jointly recover the remaining scale
coefficients. Varying $\rho$ while holding all raw scales fixed changes
baseline shares as well as substitution. Re-normalize around declared production moments.
\\[5pt]
León-Ledesma, McAdam and Willman (2010) \src{LMW10}.
Joint production/first-order-condition system with biased technical change.
&
Identification-method evidence, not a new portable value for V6 technology.
Normalized system estimation can improve separation of substitution and factor bias
relative to a single production equation.
&
Monte Carlo comparisons of CES estimators under specified data-generating processes.
Published AER article; full-text method checked in ECB working paper~1001 (2009).
&
Supports joint treatment of $\rho$, shares and technology paths.
Does not establish identification in V6, whose monopoly wedge and endogenous research
allocation introduce additional conditions.
\\[5pt]
Hirano, Kishi and Toda (2025) \src{HKT25}, Figure~1 note.
Same CES weight and factor-augmentation interpretation.
&
Illustration: $\alpha=0.5$, $\rho=2$, $A_XH=10$, $A_LL=1$, $n_0=1$.
Thus $A_X=10/H$ and $A_L=1/L$ in those units, not universal scale values $10$ and $1$.
&
Closed-economy theoretical figure with a generational period;
no fitted population or sampling uncertainty.
Its post-switch exponents are positive, unlike the V6 zero-exponent case.
&
Useful numerical comparison only. The figure does not separately discipline
US/RoW or the two constant V6 post-switch scales.
Do not reinterpret this illustration as evidence for a narrow productivity range.
\\[5pt]
\end{evidence}

\clearpage
\subsection{Explicit V6 mapping and proposed assignment}

\textbf{Our derivation from the V6 equations.}
For each country and reference observation, suppress country subscripts and set
$X_0=\varphi_0H$. Define the intermediate expenditure share, including its monopoly profits,
\[
s_{X,0}\equiv\frac{P_0X_0}{Y_0}
=\frac{w_{H,0}\varphi_0H}{\vartheta Y_0},
\qquad 1-s_{X,0}=\frac{w_{L,0}L}{Y_0}.
\]
The first share is \emph{not} total skilled payroll over output: research payroll is excluded,
while profits on current intermediate production are included through $1/\vartheta$.
Matched observations must satisfy these two shares adding to one under the maintained
production accounting. If they fail, report the scope/measurement residual rather than
treating inconsistent share estimates as independent exact restrictions.
Given $0<s_{X,0}<1$, fixed $\alpha\in(0,1)$ and $\rho\ne1$,
\begin{align}
B_X(\alpha,\rho;s_{X,0},Y_0,X_0)
&\equiv \frac{Y_0}{X_0}
 \left(\frac{s_{X,0}}{\alpha}\right)^{1/(1-\rho)}
= A_{X,0}, \label{eq:bx}\\
B_L(\alpha,\rho;s_{X,0},Y_0,L)
&\equiv \frac{Y_0}{L}
 \left(\frac{1-s_{X,0}}{1-\alpha}\right)^{1/(1-\rho)}
= A_{L,0}. \label{eq:bl}
\end{align}
These are conditional production inversions, not an independent demonstration that
research free entry and the complete equilibrium support the chosen $\varphi_0$.

\begin{evidence}
M11: $\alpha_{US},\alpha_W$.
&
\textbf{Proposed convention:} $\alpha_i=0.5$ with joint rescaling of augmentations,
when $\rho_i\ne1$. No empirical confidence interval is assigned.
&
Only $\alpha_iA_{X,i}^{1-\rho_i}$ and
$(1-\alpha_i)A_{L,i}^{1-\rho_i}$ enter the raw CES.
An alternative $\alpha_i'$ is equivalent only if both combinations are preserved.
&
This is a normalization, not a half-income-share finding.
At $\rho=1$, $\alpha$ becomes the Cobb--Douglas share and cannot be varied
as a pure scale convention. Near one, the inversion is ill-conditioned.
\\[5pt]
M13: $\bar A_{X,US,u}$ and $\bar A_{L,US,u}$.
&
\textbf{Proposed baseline:}
$B_{X,US}/N_{US,0}^{\xi_u}$ and
$B_{L,US}/N_{US,0}^{\nu_u}$.
No portable numerical level is assigned before matching data.
&
Initial US output, production-skilled and unskilled inputs, matched wages,
markup, and a declared $N_{US,0}$ normalization. These production moments
become calibration targets if used here.
&
Propagate wage/share, research-allocation and $\rho$ uncertainty jointly.
Changing units changes the bars. A range imported from a different normalization
would create artificial precision.
\\[5pt]
M15: $\bar A_{X,W}$ and $\bar A_{L,W}$.
&
\textbf{Proposed baseline:} $B_{X,W}$ and $B_{L,W}$ for the declared RoW aggregate.
No numerical full-RoW candidate is presently supported.
&
Matched country coverage, output units, skill definition and wage accounting.
CC06/Rossi provide measurement designs, not a ready-made V6 aggregate.
&
Use alternative coverage and input definitions to construct joint ranges.
Maintain the US--RoW relative scale under any common output-unit normalization.
Covered-country estimates must retain their coverage label.
\\[5pt]
M14: $\bar A_{X,US,b}$ and $\bar A_{L,US,b}$.
&
\textbf{Proposed reference scenario:}
$\bar A_{X,US,u}N_\dagger^{\xi_u}$ and
$\bar A_{L,US,u}N_\dagger^{\nu_u}$ at a declared reference switch stock $N_\dagger$.
&
A continuity assumption at one reference state, not an estimated post-switch technology.
Uninterrupted $u$ observations cannot identify these counterfactual levels.
&
Robustness: separate proportional level shocks $m_X,m_L$ around the reference.
No empirical bounds verified. Optional $m_j\in\{0.8,1,1.2\}$ is an analyst stress grid only.
Constant bars cannot impose continuity at every possible random switch date.
\\[5pt]
\end{evidence}

\textbf{Interpretation of the post-switch scenario.}
Reference continuity pertains to factor augmentations only. Output, wages and equity values
may still jump when research allocation or expectations change.
The $(m_X,m_L)$ grid is not a literature-derived plausible interval and is not recommended
as a substitute for identifying the economic transition. Its role is to expose how conclusions
depend on an unobserved technology level. Reassess the internally calibrated financial block
for each maintained external scenario.
